\section{Solution Overview}

\subsection{Design principles}

\begin{description}[leftmargin=0pt,style=nextline]
\item[Enforce, do not request.] Isolation is implemented where the application
cannot override it: kernel network namespaces and firewall rules, not
configuration flags.
\item[Configuration over code.] The model catalogue and the domain ontology are
YAML documents. Adding a model or a domain changes no Python.
\item[Embedded over networked.] The graph database and vector store run
in-process. On an air-gapped deployment there is no listening port for anything
to reach, and no service to orchestrate.
\item[Degrade, never fail.] If the embedding model is absent, graph traversal
still answers structural questions. If a tool errors, the failure is reported
to the agent rather than injected into a deliverable.
\item[Provenance everywhere.] Chunks, graph nodes and graph edges all carry
document, page and region identifiers, so any statement can be traced back.
\end{description}

\subsection{Capability map}

\begin{center}
\renewcommand{\arraystretch}{1.25}
\begin{tabularx}{\linewidth}{@{}l X l@{}}
\toprule
\textbf{Capability} & \textbf{Mechanism} & \textbf{Module} \\
\midrule
Multi-model serving & Declarative manifest, capability tags, priorities &
  \code{config/models.yaml} \\
Task routing & Two-stage: heuristics first, 4B classifier fallback &
  \code{app/router/classifier.py} \\
Memory management & Budgeted LRU residency with eviction &
  \code{app/router/manager.py} \\
Agentic execution & Plan $\rightarrow$ Execute $\rightarrow$ Critique loop &
  \code{app/agent/loop.py} \\
Tool layer & 12 schema-declared tools with argument coercion &
  \code{app/tools/} \\
Sandboxing & Unprivileged user + network namespace &
  \code{app/tools/sandbox.py} \\
Document understanding & ONNX OCR on CPU, vision model on GPU &
  \code{app/tools/vision.py} \\
Knowledge graph & Typed extraction, canonicalisation, embedded store &
  \code{app/graphrag/} \\
Hybrid retrieval & Dense + BM25 fused by reciprocal rank &
  \code{app/retrieval/vector.py} \\
Deliverables & Template-filled Office documents &
  \code{app/tools/docgen.py} \\
Sovereignty & Egress monitor, canary, hash-chained audit &
  \code{app/sentinel/} \\
\bottomrule
\end{tabularx}
\end{center}

\subsection{A worked example}

The following illustrates the full path through the system.

\begin{enumerate}
\item An engineer uploads a scanned inspection report and asks: \emph{``Which
      units are affected if P-101A is isolated, and draft an approval note
      recommending seal replacement.''}
\item The \term{router} classifies the request. The presence of a document
      attachment and deliverable vocabulary routes planning to the 14B
      reasoning model.
\item The \term{agent} produces a plan: query the graph for P-101A's
      neighbourhood; compute the impact; generate an approval note.
\item \term{graph\_query} traverses three hops from \code{Equipment:P-101A} and
      returns E-204 (one hop, direct feed), V-1201 (two hops, indirect) and
      TI-2043 (the instrument monitoring E-204).
\item The result is \term{rendered} into readable prose and substituted into
      the document-generation step through a \code{\{\{step1\}\}} placeholder.
\item \term{make\_approval\_note} writes a real \code{.docx} with headings, an
      approval signature table and a provenance footer.
\item Throughout, the \term{sentinel} panel shows
      \code{external\_connections: 0}, and the audit chain records every prompt,
      model choice and tool call.
\end{enumerate}

The critical detail is step 4. V-1201 is never mentioned in the same sentence
as P-101A anywhere in the corpus. It is reachable only by following
\code{P-101A --FEEDS--> E-204 --FEEDS--> V-1201}. No embedding model retrieves
it; a graph walk finds it every time.
